\documentclass[10pt, a4paper, twocolumn]{article}

\usepackage[utf8]{inputenc}
\usepackage[T1]{fontenc}
\usepackage{times}
\usepackage{microtype}
\usepackage{graphicx}
\usepackage{booktabs}
\usepackage{hyperref}
\usepackage{url}
\usepackage{amsmath}
\usepackage{cite}
\usepackage{geometry}
\geometry{a4paper, top=2.5cm, bottom=2.5cm, left=2cm, right=2cm, columnsep=0.6cm}

% ---------------------------------------------------------------
\title{\textbf{Frame Annotation of Legal Reporting Obligations:\\
A Dataset and Baseline}}
 
\author{
  Célestine Lecat$^{1}$ \quad Víctor Rodríguez-Doncel$^{2}$ \quad Monica Palmirani$^{1}$ \\
  \\
  $^{1}$Università di Bologna, Bologna, Italy \\
  $^{2}$Universidad Politécnica de Madrid, Madrid, Spain \\
  \\
  \texttt{\{celestine.lecat, monica.palmirani\}@unibo.it} \\
  \texttt{vrodriguez@fi.upm.es}
}

\date{}
% ---------------------------------------------------------------

\begin{document}

\maketitle

\begin{abstract}
TODO. This paper presents a gold standard dataset of legal provisions
annotated with FrameNet frames, focusing on reporting obligations
in European Union law. We define a domain-specific frame inventory,
describe the annotation methodology, and provide a baseline annotation
system built on a large language model. We discuss applications of the
resource to legal knowledge representation, compliance analysis, and
frame-based reasoning.
\end{abstract}

% ---------------------------------------------------------------
\section{Introduction}
\label{sec:intro}

Reporting obligations are among the most pervasive instruments of European
Union regulation. Across domains as diverse as banking, insurance, financial
markets, data protection, environmental protection and product safety, EU
legislation routinely requires regulated actors to transmit specified
information to a competent authority, frequently on a recurring basis and
within strict time limits. A credit institution must file prudential reports to
its supervisor at regular intervals; a data controller must notify the
supervisory authority of a personal-data breach without undue delay and, where
feasible, within seventy-two hours; an operator must report emissions or
incidents to a designated agency. Despite their diversity, such provisions
share a recognisable semantic structure: an obligated party, an act of
submitting or notifying, a recipient authority, the information to be reported,
and temporal parameters expressing deadlines, periodicity or reference points.

The volume and heterogeneity of these obligations has itself become a
regulatory concern. Scattered across hundreds of legal acts and drafted in
dense natural language, reporting requirements are costly to identify, monitor
and keep current, which has motivated efforts both to rationalise them at the
policy level and to represent them formally. On the formal side, Palmirani et
al.~\cite{palmirani-2025-reporting} propose an ontology dedicated to reporting
requirements in European legislation, modelling the actors involved, the
content to be reported and the temporal conditions that govern it. Populating
such a model, however, presupposes the ability to locate and interpret the
relevant provisions in the legal text itself, a step that remains largely
manual.

Frame semantics offers a natural intermediate representation for this step.
Prior work on contract law, for example
ContractFrames~\cite{navas-loro-2019-contractframes}, has shown that
frame-based structures can bridge natural-language legal text and formal logic
by capturing events and their participants in a normalised form. We pursue an
analogous strategy for regulatory reporting: we annotate the provisions that
impose reporting duties with FrameNet~\cite{baker-etal-1998-berkeley} frames,
making explicit who must report what, to whom, and when, and we serialise the
result as interoperable linked data so that it can feed ontologies such as the
one above.

In this paper we make the following contributions. First, we define a
domain-specific inventory of FrameNet frames tailored to reporting obligations,
adapting general-purpose frames so that their frame elements align with the
participants and temporal parameters characteristic of this domain. Second, we
present a gold-standard dataset of EU legal provisions annotated with this
inventory, together with the annotation methodology. Third, we provide a
baseline annotation system built on a large language model and evaluate it
against the gold standard. Finally, we discuss applications of the resource to
legal knowledge graphs and compliance analysis.

% ---------------------------------------------------------------
\section{State of the Art}
\label{sec:soa}

\textbf{FrameNet.}
FrameNet~\cite{baker-etal-1998-berkeley} is a lexical database of English grounded in
frame semantics~\cite{fillmore-1982-frame}, the theory that words evoke structured
conceptual frames representing stereotypical situations, events or relations.
Each frame defines a set of participants called Frame Elements (FEs), which may
be core (semantically obligatory) or non-core (peripheral). FrameNet 1.7
contains over 1{,}200 frames, 13{,}000 lexical units and 200{,}000 annotated
example sentences. Frames are interconnected through ontological relations such
as \textit{Inheritance}, \textit{Using}, \textit{Subframe} and
\textit{Perspective\_on}, forming a rich semantic network that supports both
lexicographic and computational applications.

\textbf{Framester.}
Framester~\cite{gangemi2016framester} is a large-scale RDF/OWL linked data hub
developed at CNR and Università di Bologna that integrates FrameNet with
WordNet, VerbNet, BabelNet, DBpedia, Yago and the DOLCE-Zero foundational
ontology into a single strongly connected knowledge graph. Frames are
formalised as OWL classes, frame elements as properties, and lexical units as
linked data resources, enabling full-fledged SPARQL querying and OWL reasoning
over frame-based knowledge. Framester exposes a public SPARQL
endpoint\footnote{\url{https://w3id.org/framester/sparql}} and its data is
available on GitHub.\footnote{\url{https://github.com/framester/Framester}}
In this work, Framester URIs are used to serialise frame annotations as
interoperable Turtle RDF using the NIF standard~\cite{hellmann2013nif}.

\subsection{Frame Semantic Parsing}

Frame semantic parsing is the task of automatically producing FrameNet-style
analyses of text, conventionally decomposed into three interdependent
subtasks: frame identification (selecting the frame evoked by a target),
argument identification (delimiting the spans that fill its roles) and role
classification (labelling those spans with frame elements). A representative
recent approach is the double-graph framework of
Zheng~et~al.~\cite{zheng-etal-2022-double}, whose KID model (Knowledge-guided
Incremental semantic parser with Double-graph) recasts parsing as incremental
graph construction. KID couples two graphs: a \textit{Frame Knowledge Graph},
a static heterogeneous graph that encodes frames and their frame elements
together with the ontological relations declared in FrameNet, and a
\textit{Frame Semantic Graph}, the per-sentence structure built up
incrementally from the text. By letting knowledge-enhanced frame and
frame-element representations from the former guide the construction of the
latter, the model strengthens the interactions between the three subtasks and
the dependencies among arguments that pipeline systems treat in isolation,
reporting improvements of up to 1.7~F$_1$ over the previous state of the art
on two FrameNet benchmarks. This separation between an ontological knowledge
layer and an instance-level semantic graph is conceptually close to the
representation we adopt, although our baseline delegates the parsing itself to
a large language model rather than learning a dedicated graph-construction
model (Section~\ref{sec:baseline}).

\subsection{Frame Semantics in Legal Text}

The application of frame semantics to legal language has a smaller but growing
body of work. Closest to our setting, ContractFrames~\cite{navas-loro-2019-contractframes}
applies a frame-based representation to contract law, extracting events related
to the status of a purchase contract into frames and mapping them to logic
clauses of the PROLEG legal-reasoning system through a dedicated Contract
Workflow Ontology; it thereby uses frames as an intermediate layer between
natural-language contractual text and formal logic, a role analogous to the one
we envisage for reporting obligations. More broadly, the enrichment of legal
documents with interoperable linguistic annotations has been demonstrated at
platform scale by the Lynx project~\cite{moreno-schneider-2022-lynx}, whose
service platform processes, enriches and analyses documents from the legal
domain and represents the resulting linguistic and semantic annotations as
linked data using the NIF standard. Our work follows the same linked-data
annotation philosophy, narrowing the focus to FrameNet frames for reporting
obligations and serialising them through Framester URIs.

Most directly related to our domain, Palmirani et
al.~\cite{palmirani-2025-reporting} propose a Reporting Requirements Ontology
for European legislation, a conceptual model that formalises the obligations
imposed on regulated actors to submit information to authorities, capturing the
parties involved, the content to be reported, and the temporal conditions such
as deadlines and periodicity. Whereas that work operates at the ontological
level, defining the classes and properties needed to represent reporting
duties, our contribution is complementary and operates at the linguistic level:
we annotate the textual provisions themselves with FrameNet frames, yielding
the surface evidence from which such ontological structures can be populated.
The frame inventory we adopt is in fact designed so that its frames and frame
elements align with the participants and temporal parameters of reporting
obligations, providing a natural bridge towards ontologies of this kind.

% ---------------------------------------------------------------
\section{Dataset}
\label{sec:dataset}

TODO.

\subsection{Source Corpus}

TODO.

\subsection{Annotation Methodology}

TODO.

\subsection{Dataset Statistics}

TODO.

% ---------------------------------------------------------------
\section{Baseline Annotation System}
\label{sec:baseline}

TODO.

\subsection{System Description}

TODO.

\subsection{Evaluation}

TODO.

% ---------------------------------------------------------------
\section{Applications}
\label{sec:applications}

TODO.

\subsection{Legal Knowledge Graphs}

TODO.

\subsection{Compliance Analysis}

TODO.

% ---------------------------------------------------------------
\section{Conclusion}
\label{sec:conclusion}

TODO.

% ---------------------------------------------------------------
\section*{Acknowledgements}

TODO.

% ---------------------------------------------------------------
\bibliographystyle{plain}
\bibliography{references}

\end{document}
